\subsection{NarraCrime Dataset}
\label{sec:narracrime}

NarraCrime is a diagnostic benchmark for non-interactive, evidence-grounded reasoning over detective narratives. Each instance provides a self-contained mystery and asks a model to identify the responsible person and justify the conclusion using evidence from the narrative.

NarraCrime-300 contains 300 synthetic mystery narratives across Easy, Medium, and Complex difficulty levels, with 100 cases per level. The cases were produced through a structured, AI-assisted workflow covering blueprint design, evidence-chain specification, distractor planning, narrative realization, and consistency checks. Difficulty is controlled through factors such as the number of suspects, required evidence cues, misleading trails, and implicit-premise dependencies.

Each released case contains a narrative, a reference answer, predefined evidence cues, and a structured annotation. The annotation records the culprit, intent, action schema, supporting evidence, distractors, and a construction blueprint. These fields support verdict, content-recovery, and evidence-grounding evaluation.

The task and evaluation design were informed by the non-interactive detective-reasoning setting explored in SABA. In particular, NarraCrime's Intent Recall, Action Schema Recall, and Evidence Coverage adapt SABA's motive-recall, modus-operandi-recall, and clue-coverage dimensions to the NarraCrime annotation protocol. EVAR additionally reports role-aware verdict scoring and final-claim reliability measures. We therefore do not present all evaluation dimensions as originating with EVAR.
